% SOLUCIÓN Ejercicio 3 — Informe: amplificador operacional inversor
\documentclass[12pt,a4paper]{article}
\usepackage[utf8]{inputenc}
\usepackage[T1]{fontenc}
\usepackage[spanish]{babel}
\usepackage{amsmath,amssymb}
\usepackage{booktabs}
\usepackage{siunitx}
\sisetup{locale=FR, output-decimal-marker={,}}
\usepackage{circuitikz}
\usepackage[hidelinks]{hyperref}
\usepackage{float}

\begin{document}
\begin{titlepage}
  \centering
  {\Large Universidad Nacional Tecnológica de Lima Sur\par}
  {\large Facultad de Ingeniería Electrónica\par}
  \vspace{2cm}
  {\Huge\bfseries Circuito Amplificador Operacional Inversor\par}
  \vspace{1cm}
  {\large Informe de Laboratorio N.° 3\par}
  \vfill
  {\large Estudiante UNTELS\par}
  {\large \today\par}
\end{titlepage}
\tableofcontents
\newpage

\section{Objetivos}
\begin{itemize}
  \item Verificar la ganancia en lazo cerrado del amplificador inversor.
  \item Comparar la salida medida con la teórica.
\end{itemize}

\section{Marco teórico}
Con un amplificador operacional ideal, la tierra virtual en la entrada inversora conduce a
\begin{equation}
  A_v = \frac{V_{out}}{V_{in}} = -\frac{R_f}{R_1}
  \label{eq:ganancia}
\end{equation}
y la salida queda limitada por la alimentación:
\begin{equation}
  |V_{out}| \le V_{sat} \approx V_{CC} - \SI{1,5}{\volt}
  \label{eq:sat}
\end{equation}

\section{Circuito}
La Figura~\ref{fig:inversor} muestra el montaje con $R_1=\SI{10}{\kilo\ohm}$ y $R_f=\SI{100}{\kilo\ohm}$.
\begin{figure}[htbp]
  \centering
  \begin{circuitikz}[american]
    \draw (0,0) node[op amp] (A) {};
    \draw (A.-) -- ++(-0.6,0) coordinate (fb);
    \draw (fb) to[R=$R_1$, -o] ++(-2.2,0) node[left] {$V_{in}$};
    \draw (A.+) -- ++(-0.4,0) node[ground] {};
    \draw (fb) to[short, *-] ++(0,1.7) coordinate (top);
    \draw (top) to[R=$R_f$] (top -| A.out) -- (A.out);
    \draw (A.out) to[short, *-o] ++(1,0) node[right] {$V_{out}$};
  \end{circuitikz}
  \caption{Amplificador operacional inversor}
  \label{fig:inversor}
\end{figure}

\section{Resultados}
Por la ecuación~\eqref{eq:ganancia}, $A_v=-10$. La Tabla~\ref{tab:res} contiene valores
\emph{simulados} (reemplace por sus mediciones).
\begin{table}[htbp]
  \centering
  \caption{Tensión de salida del inversor ($A_v=-10$)}
  \label{tab:res}
  \begin{tabular}{S[table-format=1.2] S[table-format=-2.2] S[table-format=-2.2]}
    \toprule
    {$V_{in}$ (\si{\volt})} & {$V_{out}$ teórico (\si{\volt})} & {$V_{out}$ medido (\si{\volt})} \\
    \midrule
    0,10 & -1,00 & -0,98 \\
    0,20 & -2,00 & -1,97 \\
    0,50 & -5,00 & -4,95 \\
    0,80 & -8,00 & -7,93 \\
    1,00 & -10,00 & -9,90 \\
    \bottomrule
  \end{tabular}
\end{table}

\section{Conclusiones}
La salida medida no difiere más del 2\,\% de la teórica; la diferencia se atribuye a la
tolerancia de las resistencias y a la ganancia finita del operacional.

\begin{thebibliography}{9}
  \bibitem{sedra} Sedra, A. S., \& Smith, K. C. (2020). \emph{Microelectronic circuits} (8.\textsuperscript{a} ed.). Oxford University Press.
  \bibitem{horowitz} Horowitz, P., \& Hill, W. (2015). \emph{The art of electronics} (3.\textsuperscript{a} ed.). Cambridge University Press.
\end{thebibliography}
\end{document}
